\section{Known theoretical and numerical issues}
\label{sec:known-issues}

\subsection{Internal control}

\begin{remark}[Ill-posedness of the unpenalized minimization]
  \label{rem:coercivity-space}
  The space on which $J$ (Proposition~\ref{prop:pde-gramian}) is coercive is
  much larger than $L^2(0,1)$: it is the completion of $L^2(0,1)$ with respect
  to the norm
  $\big(\tfrac12\int_0^1\!\int_0^T|\mathbf{1}_\omega\varphi|^2\,dx\,dt\big)^{1/2}$
  -- exactly the space $F$ of Remark~\ref{rem:hum-existence}. This space can
  hardly be approximated by a finite-dimensional basis
  \citep{munch2010numerical}, which is why the numerical problem of minimizing
  the dual functional is intrinsically ill-posed.
\end{remark}

Two genuine cost-of-control results appear in the source material: an
exponential-in-\emph{time} estimate,
$\|v\|_{L^\infty(\Sigma)} \le e^{-2\lambda_1 T}C(T-\tau)\|z_0\|_{L^2(\Omega)}$
(here $\lambda_1$ is the first Dirichlet-Laplacian eigenvalue on $\Omega$,
$z_0$ the shifted initial datum, and $\tau$ an intermediate time -- all from
the proof of Theorem~\ref{thm:constrained-linear} in
Section~\ref{sec:constrained-aside}, cost in $T$, not in $\varepsilon$), and a
purely \emph{experimental} control-norm blow-up as the control region $\omega$
shrinks. No exponential-in-$\varepsilon$ blow-up result
(of the form $k\sim e^{1/\varepsilon}$, sometimes seen in the wider numerical
HUM literature) is established in the source material used here, and none is
claimed.

\begin{remark}[Gramian conditioning vs.\ control-region geometry]
  \label{rem:gramian-geometry}
  Although the literature guarantees exact/null controllability for
  internal control on any open $\omega$, with no restriction on its geometry
  \citep{zuazua2002controllability}, the condition number of the Gramian
  $\Lambda$ depends strongly on the location of $\omega$; when discretized in
  space, $\Lambda_N$, it blows up as $N$ increases and $\omega$ shrinks
  \citep{munch2011optimal}.
\end{remark}

\begin{remark}[Time-integrator sensitivity]
  Convergence of the numerical HUM method is highly sensitive to the choice
  of time integrator; explicit schemes blow up for most mesh-parameter pairs,
  even when the underlying finite-difference scheme is convergent and
  uniformly controllable (Section~\ref{sec:discrete}). Gradient-type
  optimization methods can likewise fail even under those same favorable
  conditions; preconditioning or other iterative methods for the resulting
  large linear systems are a natural, and currently active, avenue for
  improving conjugate-gradient convergence (Theorem~\ref{thm:cg-rate}).
\end{remark}

\begin{remark}[Regularity of the adjoint state and the penalization norm]
  There is a strong dependency between the regularity of the adjoint state
  and the minimizer of the dual functional: different choices of penalization
  norm produce genuinely different solution behavior
  \citep{kindermann1999convergence}.
\end{remark}

\subsection{Boundary control}

Boundary control requires a change of norms and function spaces relative to
internal control (Remark~\ref{rem:boundary-h10}), which in practice makes the
numerical method more unstable. For a uniform control-point distribution, a
denser distribution near the endpoints of the control interval may improve
results \citep{kalimeris2021numerical}.

\begin{remark}[Lack of robustness]
  Small perturbations in input data can produce widely different numerical
  results -- a general robustness concern for both control types, not
  specific to either.
\end{remark}
